\chapter{Systemic Autoimmune and Connective Tissue Diseases}
\label{ch:connective_tissue}
Systemic autoimmune and connective tissue diseases encompass a heterogeneous group of disorders characterized by immune dysregulation, autoantibody production, and chronic inflammation affecting multiple organ systems. These conditions include systemic connective tissue diseases (scleroderma, polymyalgia rheumatica, mixed connective tissue disease, relapsing polychondritis, IgG4-related disease, eosinophilic fasciitis), systemic vasculitides affecting vessels of all sizes, and other systemic autoimmune disorders such as adult-onset Still's disease, sarcoidosis, amyloidosis, hemophagocytic lymphohistiocytosis, and periodic fever syndromes. Many share overlapping clinical and serologic features, and a multidisciplinary approach is essential for diagnosis and management.

%-------------------------------------------------------------
\section{Other Systemic Autoimmune Disorders}
%-------------------------------------------------------------

%-------------------------------------------------------------

%-------------------------------------------------------------

%-------------------------------------------------------------

%-------------------------------------------------------------

%-------------------------------------------------------------

%-------------------------------------------------------------

\label{sec:Other Systemic Autoimmune Disorders}
%-------------------------------------------------------------%-------------------------------------------------------------

\subsection{Adult-Onset Still's Disease}
\label{sec:Adult-Onset Still's Disease (AOSD)}
Adult-onset Still's disease is a rare systemic inflammatory disorder of unknown etiology, considered the adult counterpart of systemic juvenile idiopathic arthritis (sJIA), characterized by a classic triad of high-spiking fevers, evanescent salmon-pink rash, and arthralgia/arthritis, with an incidence of 0.16--0.4 per 100,000, a bimodal age distribution (15--25 and 36--46 years), and a slight female predominance. The condition results from dysregulated innate immunity with activation of the inflammasome, leading to overproduction of IL-1, IL-6, IL-18, and other pro-inflammatory cytokines, with macrophage activation being a key feature. Clinical features include quotidian or double-quotidian fevers, an evanescent, salmon-pink, non-pruritic maculopapular rash (Koebner phenomenon), arthralgia and arthritis (initially pauciarticular, may progress to severe destructive polyarthritis, particularly involving wrists and carpal bones---carpal ankylosis is characteristic), sore throat (intense, exudative, a highly distinctive feature), lymphadenopathy and splenomegaly, serositis, and myalgia. Diagnosis is by Yamaguchi criteria, with serum ferritin characteristically markedly elevated (often $>$5--10 times normal, $>$1000 ng/mL), with decreased glycosylated ferritin ($<$20\%). Management is stratified: NSAIDs for mild disease, corticosteroids for moderate disease, and IL-1 inhibitors (anakinra, canakinumab) or IL-6 inhibitors (tocilizumab) for refractory or severe disease. Prognosis is variable; some patients achieve remission, while others have a chronic relapsing course.

\subsection{Amyloidosis}
\label{sec:Amyloidosis}
Amyloidosis is a group of diseases characterized by the extracellular deposition of abnormal, insoluble fibrillar proteins (amyloid) in various tissues and organs, disrupting normal structure and function, with three main types clinically important: AL amyloidosis (immunoglobulin light chain---the most common systemic form, associated with plasma cell dyscrasia/multiple myeloma), AA amyloidosis (secondary/reactive; serum amyloid A protein derived from chronic inflammation), and ATTR amyloidosis (transthyretin; hereditary or wild-type/senile systemic amyloidosis). The incidence of AL amyloidosis is 3--5 per million, typically presenting at age 60--70. Clinical features depend on the organs involved: AL amyloidosis commonly presents with nephrotic syndrome (renal amyloidosis---most common presentation, 70\%), restrictive cardiomyopathy (heart failure with preserved ejection fraction, low-voltage QRS on ECG, thickened ventricular walls on echocardiography with a ``speckled'' appearance), hepatomegaly, peripheral and autonomic neuropathy, macroglossia (tongue enlargement---pathognomonic for AL), periorbital purpura, and carpal tunnel syndrome. Diagnosis requires tissue biopsy with Congo red staining and apple-green birefringence under polarized light---the diagnostic gold standard. Management of AL amyloidosis is anti-plasma cell therapy (similar to multiple myeloma: bortezomib-based regimens, autologous stem cell transplantation in selected patients); AA amyloidosis is treated by controlling the underlying inflammatory condition; ATTR amyloidosis is treated with tafamidis (stabilizes transthyretin, slows neuropathy and cardiomyopathy progression). Prognosis is poor for AL amyloidosis without treatment but has improved significantly with modern therapies.

\subsection{Hemophagocytic Lymphohistiocytosis and Macrophage Activation Syndrome}
\label{sec:Hemophagocytic Lymphohistiocytosis (HLH)}
Hemophagocytic lymphohistiocytosis is a life-threatening syndrome of pathologic immune activation, characterized by uncontrolled proliferation of activated lymphocytes and macrophages (histiocytes) producing a massive cytokine storm (hypercytokinemia: IFN-$\gamma$, IL-1, IL-6, IL-18, TNF). There are two forms: primary (familial) HLH (fHLH), caused by autosomal recessive mutations in genes critical for perforin-mediated cytotoxic function, typically presenting in infancy or early childhood; and secondary HLH (sHLH), triggered by infections (EBV is the most common trigger), malignancies, autoimmune diseases (macrophage activation syndrome, MAS, associated with systemic juvenile idiopathic arthritis [sJIA] and adult-onset Still's disease [AOSD]), and immunosuppression. The incidence of HLH is estimated at 1 per 50,000--100,000 in children. Clinical features include persistent high fever, hepatosplenomegaly, lymphadenopathy, cytopenias (at least two lineages: anemia, thrombocytopenia, neutropenia), skin rash, and neurological symptoms (seizures, altered mental status, meningismus). Diagnosis is by HLH-2004 criteria (five of eight: fever, splenomegaly, cytopenias, hypertriglyceridemia and/or hypofibrinogenemia, hemophagocytosis in bone marrow/spleen/lymph nodes, low or absent NK-cell activity, ferritin $\geq$500 ng/mL, elevated soluble CD25/IL-2R $\geq$2400 U/mL). Management of HLH includes immunosuppression with dexamethasone and etoposide (HLH-94/2004 protocol), with allogeneic hematopoietic stem cell transplantation for familial HLH and refractory cases. Prognosis is poor without treatment but has improved with protocolized therapy.

\subsection{Periodic Fever Syndromes}
\label{sec:Periodic Fever Syndromes}
Periodic fever syndromes (also called autoinflammatory diseases) are a group of inherited disorders of the innate immune system characterized by recurrent, self-limited episodes of fever and systemic inflammation in the absence of autoantibodies or antigen-specific T cells. They are distinct from autoimmune diseases and arise from dysregulation of the inflammasome and IL-1$\beta$ production (the ``IL-1opathies'').

\subsubsection{Familial Mediterranean Fever}
Familial Mediterranean fever is the most common hereditary periodic fever syndrome, caused by mutations in the \textit{MEFV} gene encoding pyrin, a component of the pyrin inflammasome, most prevalent in Mediterranean populations (Armenian, Turkish, Arab, Jewish---carrier frequency 1:5 in Armenians). Clinical features begin in childhood (90\% before age 20) and consist of recurrent, self-limited attacks (12--72 hours) of high fever, serositis (peritonitis---severe abdominal pain resembling acute abdomen, pleuritis, pericarditis), arthritis (usually monoarticular of the lower extremity), and an erysipelas-like erythema on the lower leg, with the diagnostic hallmark being a dramatic response to colchicine. The most serious long-term complication is amyloidosis (AA type) if untreated. Diagnosis is clinical (Tel Hashomer criteria) or by genetic confirmation of biallelic \textit{MEFV} mutations. Management is daily colchicine (0.5--2 mg/day), which prevents attacks and amyloidosis in 95\% of patients, with IL-1 inhibitors (anakinra, canakinumab) used for colchicine-resistant FMF. Prognosis is excellent with colchicine adherence; AA amyloidosis risk is nearly eliminated.

\subsubsection{TNF Receptor-Associated Periodic Syndrome}
TNF receptor-associated periodic syndrome is an autosomal dominant periodic fever syndrome caused by mutations in the \textit{TNFRSF1A} gene encoding the p55 TNF receptor, with approximately 1000 reported cases worldwide, primarily in Northern European populations. Clinical features include prolonged attacks (1--3 weeks, longer than FMF) of fever, severe myalgia (migratory with overlying erythema, a hallmark), arthralgia, abdominal pain, painful conjunctivitis and periorbital edema, erythematous rash, and serositis, with amyloidosis occurring in 10--25\% of untreated patients. Diagnosis is clinical with genetic confirmation of \textit{TNFRSF1A} mutations. Management includes corticosteroids (effective for acute attacks), IL-1 inhibitors (anakinra, canakinumab---first-line for severe disease), and etanercept (TNF receptor fusion protein). Prognosis is variable, with amyloidosis risk reduced with effective anti-IL-1 therapy.

\subsubsection{Cryopyrin-Associated Periodic Syndromes}
Cryopyrin-associated periodic syndromes are a spectrum of autoinflammatory diseases caused by gain-of-function mutations in the \textit{NLRP3} gene encoding cryopyrin (NALP3), a key component of the NLRP3 inflammasome, encompassing three overlapping phenotypes in order of increasing severity: familial cold autoinflammatory syndrome (FCAS---mildest), Muckle-Wells syndrome (MWS---intermediate), and neonatal-onset multisystem inflammatory disease/chronic infantile neurologic cutaneous articular syndrome (NOMID/CINCA---most severe), with an incidence of approximately 1 per 1,000,000. Clinical features include recurrent fever, urticaria-like rash (non-pruritic, cold-induced in FCAS, persistent in MWS/NOMID), arthralgia or arthritis, conjunctivitis/episcleritis, progressive sensorineural hearing loss (in MWS and NOMID), and neurological involvement (aseptic meningitis in NOMID). Diagnosis is clinical with genetic confirmation of \textit{NLRP3} mutations. Management is with IL-1 inhibitors: anakinra, canakinumab (long-acting, FDA-approved for CAPS), and rilonacept (IL-1 trap), which are dramatically effective, normalizing symptoms and inflammatory markers within hours and preventing AA amyloidosis and hearing loss progression.

\subsubsection{Hyper-IgD Syndrome}
Hyper-IgD syndrome (also known as mevalonate kinase deficiency, MKD) is an autosomal recessive periodic fever syndrome caused by mutations in the \textit{MVK} gene encoding mevalonate kinase, an enzyme in the cholesterol biosynthesis pathway, leading to accumulation of mevalonic acid and activation of the inflammasome, with approximately 300 reported cases, predominantly in Northern European populations (especially Dutch). Clinical features begin in infancy and consist of recurrent attacks (3--7 days) of high fever, lymphadenopathy (cervical), abdominal pain, vomiting, diarrhea, arthralgia/arthritis, oral ulcers, and an erythematous macular or urticarial rash, with attacks triggered by vaccinations, infections, stress, or trauma. Diagnosis is confirmed by genetic testing (biallelic \textit{MVK} mutations), with serum IgD elevated (often $>$100 IU/mL) in 80--90\% of patients. Management is supportive during attacks (NSAIDs, corticosteroids), with IL-1 inhibitors and TNF inhibitors used for severe disease. Prognosis is good; attacks decrease in frequency and severity with age, and life expectancy is normal.

\subsection{Sarcoidosis}
\label{sec:Sarcoidosis (Systemic)}
Sarcoidosis is a multisystem granulomatous disease of unknown etiology, characterized by non-caseating granulomas in affected tissues. It is discussed in detail in Chapter~\ref{ch:lungs} (Section~\ref{sec:Sarcoidosis}). Briefly, it most commonly affects the lungs and intrathoracic lymph nodes ($>$90\%), but can involve any organ, including the skin (lupus pernio, erythema nodosum), eyes (uveitis, dry eyes), liver, spleen, heart (arrhythmias, cardiomyopathy, heart failure), central nervous system (neurosarcoidosis---cranial nerve palsies, meningitis, hypothalamic-pituitary involvement), exocrine glands, and bones. The diagnosis requires compatible clinical/radiographic findings, histologic demonstration of non-caseating granulomas, and exclusion of other causes. Management is indicated for organ-threatening or symptomatic disease: corticosteroids are first-line; hydroxychloroquine is used for cutaneous sarcoidosis; methotrexate, leflunomide, mycophenolate, azathioprine, and TNF inhibitors are used for refractory or severe disease. Prognosis is favorable: two-thirds of patients experience spontaneous remission within a decade.

%-------------------------------------------------------------
\section{Small Vessel Vasculitis}
%-------------------------------------------------------------

%-------------------------------------------------------------

%-------------------------------------------------------------

%-------------------------------------------------------------

%-------------------------------------------------------------

%-------------------------------------------------------------

%-------------------------------------------------------------

\label{sec:Small Vessel Vasculitis}
%-------------------------------------------------------------%-------------------------------------------------------------

\subsection{Anti-GBM Disease}
\label{sec:Anti-GBM Disease (Goodpasture Syndrome)}
Anti-glomerular basement membrane disease, also known as Goodpasture syndrome, is a rare, life-threatening small-vessel vasculitis characterized by autoantibodies directed against the non-collagenous domain of the $\alpha$3 chain of type IV collagen ($\alpha$3(IV)NC1), found in the basement membranes of glomeruli and pulmonary alveoli, with an incidence of 0.5--1 per million and a bimodal age distribution (young adults 20--30 years with male predominance, and older adults 60--70 years with female predominance). The condition involves anti-GBM antibodies binding to the target antigen in the GBM and alveolar basement membrane, activating complement and recruiting neutrophils, leading to crescentic glomerulonephritis and alveolar hemorrhage, with genetic susceptibility strongly associated with HLA-DRB1*1501. Clinical features include rapidly progressive glomerulonephritis (hematuria, proteinuria, oliguria, rising creatinine, crescentic GN on biopsy), with or without pulmonary involvement (hemoptysis, dyspnea, alveolar hemorrhage---seen in 60--80\%, more common in smokers and young patients). Diagnosis is based on detection of circulating anti-GBM antibodies by ELISA and kidney biopsy showing linear IgG deposition along the GBM on immunofluorescence (the pathognomonic finding). Management is urgent: high-dose corticosteroids combined with cyclophosphamide and plasmapheresis (daily for 7--14 days to remove pathogenic antibodies). Prognosis depends on the degree of renal impairment at presentation; if dialysis-dependent at diagnosis, renal recovery is rare.

\subsection{Cryoglobulinemic Vasculitis}
\label{sec:Cryoglobulinemic Vasculitis}
Cryoglobulinemic vasculitis is a small-vessel vasculitis caused by immune complexes of cryoglobulins (immunoglobulins that precipitate at temperatures $<$37$^\circ$C and dissolve on rewarming) deposited in vessel walls, classified into three types: type I (monoclonal, associated with lymphoproliferative disorders), type II (mixed monoclonal and polyclonal, most commonly associated with hepatitis C virus---HCV), and type III (polyclonal, associated with autoimmune diseases such as SLE or Sjögren's syndrome). Mixed cryoglobulinemia (types II and III) represents 80\% of cases and is predominantly HCV-associated. Clinical features include palpable purpura (the most common feature, recurrent, involving lower extremities), arthralgias, peripheral neuropathy (symmetric, often painful), membranoproliferative glomerulonephritis (proteinuria, hematuria, hypertension, progressive renal dysfunction), and less commonly, Raynaud's phenomenon, skin ulcers, and digital gangrene. Diagnosis requires detection of serum cryoglobulins (blood collected and transported at 37$^\circ$C), low complement levels (C4 markedly decreased), positive rheumatoid factor, and evidence of HCV infection. Management depends on severity: mild disease is managed with NSAIDs, colchicine, or low-dose corticosteroids; moderate-to-severe disease requires rituximab, often combined with antiviral therapy (direct-acting antivirals for HCV) and plasmapheresis for life-threatening disease. Prognosis depends on the underlying cause and organ involvement.

\subsection{Eosinophilic Granulomatosis with Polyangiitis}
\label{sec:Eosinophilic Granulomatosis with Polyangiitis (EGPA)}
Eosinophilic granulomatosis with polyangiitis is a rare ANCA-associated vasculitis characterized by asthma, eosinophilia, and necrotizing vasculitis of small and medium-sized vessels with extravascular eosinophilic granulomas, with an incidence of 1--3 per million, equal sex distribution, and typically presenting in the fourth to fifth decade. The condition involves a Th2-driven immune response with prominent eosinophil activation and degranulation, ANCA (typically MPO-ANCA, positive in 30--50\%, more commonly in patients with renal and neurologic involvement), and eosinophilic infiltration of tissues. Clinical features evolve through three phases: prodromal (asthma, allergic rhinitis, often years before systemic disease), eosinophilic (peripheral eosinophilia $>$1500/$\mu$L, eosinophilic pneumonia, gastroenteritis), and vasculitic (systemic vasculitis with mononeuritis multiplex---the most common vasculitic manifestation, purpura, glomerulonephritis, cardiomyopathy, and skin nodules), with cardiac involvement (eosinophilic myocarditis, coronary arteritis, heart failure) being the leading cause of death. Diagnosis is based on the 2022 ACR/EULAR criteria. Management is stratified: patients without poor prognostic factors (FFS five-factor score=0) receive corticosteroids alone; patients with poor prognostic factors (FFS $\geq$1) receive corticosteroids plus cyclophosphamide or rituximab. Prognosis depends on cardiac involvement and renal function; 5-year survival exceeds 90\%.

\subsection{Granulomatosis with Polyangiitis}
\label{sec:Granulomatosis with Polyangiitis (GPA)}
Granulomatosis with polyangiitis is an ANCA-associated necrotizing granulomatous vasculitis affecting small and medium-sized vessels, with a predilection for the respiratory tract and kidneys, with an incidence of approximately 2--12 per million, equal sex distribution, and typically presenting in the fourth to sixth decade. The condition involves c-ANCA (cytoplasmic) directed against proteinase 3 (PR3), leading to neutrophil activation, degranulation, and respiratory burst, causing endothelial injury and necrotizing inflammation. The classic triad involves the upper respiratory tract (sinusitis, epistaxis, nasal crusting, saddle nose deformity, subglottic stenosis), lower respiratory tract (cough, hemoptysis, pulmonary nodules and cavitary lesions, alveolar hemorrhage), and kidney (rapidly progressive glomerulonephritis with crescents---pauci-immune on immunofluorescence), with other features including ocular involvement (scleritis, episcleritis, orbital pseudotumor), cutaneous purpura, mononeuritis multiplex, and arthralgias. Diagnosis requires histologic confirmation (nasal, endobronchial, or renal biopsy showing granulomatous inflammation, multinucleated giant cells, and necrotizing vasculitis) or serologic evidence of PR3-ANCA (positive in $>$90\% of active generalized GPA). Management is with high-dose corticosteroids combined with rituximab or cyclophosphamide for induction of remission, with maintenance therapy using rituximab (scheduled infusions every 6 months) or methotrexate/azathioprine. Prognosis has improved dramatically with modern therapy, with 5-year survival exceeding 80\%.

\subsection{IgA Vasculitis}
\label{sec:IgA Vasculitis (Henoch-Schönlein Purpura)}
IgA vasculitis is a small-vessel vasculitis characterized by IgA-dominant immune complex deposition in vessel walls, affecting the skin, joints, gastrointestinal tract, and kidneys, and is the most common vasculitis in children, with an incidence of 3--26 per 100,000 in children versus 0.1--1.8 per 100,000 in adults, with a median age at presentation of 4--6 years in children. The condition involves IgA1 deposition in small vessels with complement activation and leukocytoclasia, with triggers including upper respiratory tract infections. Clinical features include palpable purpura (non-thrombocytopenic, symmetrically distributed over the lower extremities and buttocks---the universal feature), arthritis/arthralgias (oligoarticular, non-erosive, typically of lower extremity joints), gastrointestinal involvement (colicky abdominal pain, nausea, vomiting, intussusception, gastrointestinal bleeding), and renal involvement (IgA nephropathy---hematuria, proteinuria, hypertension, more severe in adults). Diagnosis is clinical, supported by skin or renal biopsy showing leukocytoclastic vasculitis with predominant IgA deposition on immunofluorescence. Management is supportive for mild cases; corticosteroids are used for significant gastrointestinal involvement, orchitis, or severe arthritis. Prognosis is excellent in children (complete recovery in 95\%), with relapses in 30\%.

\subsection{Microscopic Polyangiitis}
\label{sec:Microscopic Polyangiitis (MPA)}
Microscopic polyangiitis is an ANCA-associated necrotizing vasculitis affecting small vessels (capillaries, venules, arterioles), with a predilection for the kidneys and lungs, but without granuloma formation (distinguishing it from GPA), with an incidence of 2--6 per million, a slight male predominance, and typically presenting in the sixth to seventh decade. The condition involves p-ANCA (perinuclear) directed against myeloperoxidase (MPO) in approximately 60--70\% of patients (PR3-ANCA in 20--30\%), with ANCA-induced neutrophil activation leading to necrotizing inflammation of small vessels. Clinical features include rapidly progressive glomerulonephritis (crescentic, pauci-immune: the most common and serious manifestation, present in $>$80\%), pulmonary capillaritis causing diffuse alveolar hemorrhage (hemoptysis, dyspnea, bilateral infiltrates), cutaneous purpura, mononeuritis multiplex, arthralgias, and myalgias. Diagnosis is based on the 2022 ACR/EULAR classification criteria, which emphasize MPO-ANCA positivity, renal involvement, and absence of granulomatous features. Management follows the same principles as GPA: high-dose corticosteroids plus rituximab or cyclophosphamide for induction, followed by maintenance therapy with rituximab, azathioprine, or methotrexate. Prognosis is similar to GPA, with 5-year survival exceeding 75\%.

%-------------------------------------------------------------
\section{Systemic Connective Tissue Diseases}
%-------------------------------------------------------------

%-------------------------------------------------------------

%-------------------------------------------------------------

%-------------------------------------------------------------

%-------------------------------------------------------------

%-------------------------------------------------------------

%-------------------------------------------------------------

\label{sec:Systemic Connective Tissue Diseases}
%-------------------------------------------------------------%-------------------------------------------------------------

\subsection{Ehlers-Danlos Syndrome}
\label{sec:Ehlers-Danlos Syndrome}
Ehlers-Danlos syndrome (EDS) is a heterogeneous group of connective tissue disorders caused by defects in collagen synthesis and processing, with a combined prevalence of 1 in 5,000. The most common subtype, hypermobile EDS (hEDS, 90\%), has unknown genetic basis. Classical EDS (cEDS) results from COL5A1/A2 mutations causing abnormal type V collagen; vascular EDS (vEDS) results from COL3A1 mutations causing abnormal type III collagen. Clinical features vary by subtype: hEDS presents with joint hypermobility, skin hyperextensibility, easy bruising, chronic pain, and recurrent dislocations; cEDS features velvety, hyperextensible skin, atrophic scars, and easy bruising; vEDS presents with thin, translucent skin, characteristic facial appearance (prominent eyes, thin nose, lobeless ears), and life-threatening arterial, intestinal, or uterine rupture. Diagnosis is based on Villefranche classification and genetic testing. Management involves symptomatic treatment: physical therapy for joint instability, pain management, wound care, and cardiovascular surveillance in vEDS. Prognosis is normal for hEDS and cEDS; vEDS has a median survival of 50 years, primarily due to arterial rupture.

\subsection{Eosinophilic Fasciitis (Shulman Syndrome)}
\label{sec:Eosinophilic Fasciitis (Shulman Syndrome)}
Eosinophilic fasciitis is a rare, fibrosing disorder of the fascia characterized by symmetrical, painful swelling and induration of the skin and subcutaneous tissue, with peripheral eosinophilia, with an incidence of approximately 1 per million, equal sex distribution, and typically affecting adults 40--60 years of age. The etiology is unknown, but triggers include vigorous exercise or trauma (in up to 50\%), infections, and medications (statins, antitubercular drugs), with pathophysiology involving inflammation and thickening of the deep fascia, infiltration by lymphocytes, eosinophils, macrophages, and plasma cells, leading to collagen deposition and fibrosis. Clinical features include acute or subacute onset of symmetrical, painful edema and induration of the arms and legs (sparing the hands, feet, and face---a distinguishing feature from systemic sclerosis), with a characteristic ``peau d'orange'' or cobblestone appearance, and the groove sign (linear depression overlying superficial veins) being a pathognomonic finding, with joint contractures developing due to fascial fibrosis. Diagnosis requires a deep incisional biopsy including skin, subcutaneous tissue, fascia, and muscle, showing fasciitis with a mixed inflammatory infiltrate and often eosinophils, with peripheral eosinophilia ($>$5\% or absolute count $>$500/$\mu$L) present in 70--90\%. Management is with corticosteroids (prednisone 0.5--1 mg/kg/day), with methotrexate, mycophenolate, and azathioprine used as steroid-sparing agents. Prognosis is variable; most patients improve with corticosteroids, but joint contractures may persist.

\subsection{IgG4-Related Disease}
\label{sec:IgG4-Related Disease}
IgG4-related disease is a fibroinflammatory condition characterized by tumefactive (mass-like) lesions, lymphoplasmacytic infiltration rich in IgG4-positive plasma cells, storiform fibrosis, and often elevated serum IgG4 levels, affecting virtually any organ with the most common sites being the pancreas (autoimmune pancreatitis type 1), salivary glands (sialadenitis, Mikulicz disease), orbits (dacryoadenitis, orbital pseudotumor), retroperitoneum (retroperitoneal fibrosis), thyroid (Riedel thyroiditis), kidneys (tubulointerstitial nephritis), lymph nodes, lungs, and biliary tree, with an incidence of 1--5 per 100,000, a male predominance (3:1), and typically presenting in the sixth decade. The condition results from a complex interplay of genetic susceptibility, adaptive and innate immune activation, with a Th2-skewed immune response, T follicular helper cells, and overproduction of IL-4, IL-5, IL-13, and TGF-$\beta$, driving IgG4 class switching, eosinophilia, and fibrosis. Clinical features are organ-dependent but commonly include painless swelling of affected organs, obstructive symptoms (jaundice from pancreatic/biliary involvement, hydronephrosis from retroperitoneal fibrosis), and constitutional symptoms (fatigue, weight loss, low-grade fever). Diagnosis requires histologic confirmation: characteristic histopathology (dense lymphoplasmacytic infiltrate with $>$10 IgG4+ plasma cells per high-power field, IgG4+/IgG+ ratio $>$40\%, storiform fibrosis, obliterative phlebitis). Management is with corticosteroids (first-line, highly effective), with rituximab, azathioprine, mycophenolate, and methotrexate used for steroid-refractory or relapsing disease. Prognosis is generally good but relapses are common (30--50\%) during steroid tapering.

\subsection{Marfan Syndrome}
\label{sec:Marfan Syndrome}
Marfan syndrome is an autosomal dominant connective tissue disorder caused by mutations in the FBN1 gene encoding fibrillin-1, affecting 1 in 5,000 individuals. The condition results from abnormal microfibril assembly, leading to reduced structural integrity of connective tissue and dysregulation of TGF-beta signaling. Clinical features include skeletal manifestations (tall stature, arachnodactyly, pectus excavatum or carinatum, scoliosis, joint hypermobility), ocular manifestations (ectopia lentis, myopia, retinal detachment), and cardiovascular manifestations (aortic root dilation, aortic aneurysm, aortic dissection, mitral valve prolapse). Diagnosis is based on the revised Ghent criteria, including systemic score, aortic root diameter Z-score, ectopia lentis, and FBN1 genetic testing. Management involves beta-blockers (atenolol) or losartan to reduce aortic growth rate, serial echocardiography for aortic surveillance, elective aortic root replacement when diameter reaches 50 mm (or earlier with family history of dissection), and lifestyle modification (avoidance of contact sports, isometric exercise). Prognosis has improved dramatically; life expectancy now approaches 70 years with appropriate surveillance and prophylactic surgery, compared to 45 years in the 1970s.

\subsection{Mixed Connective Tissue Disease}
\label{sec:Mixed Connective Tissue Disease}
Mixed connective tissue disease is a distinct systemic autoimmune disorder characterized by overlapping clinical features of systemic lupus erythematosus (SLE), systemic sclerosis (SSc), and polymyositis/dermatomyositis (PM/DM), in the presence of high-titer anti-U1 ribonucleoprotein (anti-RNP) antibodies, with a prevalence of approximately 1--2 per 100,000, a 9:1 female predominance, and typically presenting in the third to fourth decade. The condition results from autoantibodies against U1-RNP, leading to immune complex formation, complement activation, and tissue injury, with prominent endothelial and microvascular involvement. Clinical features evolve over time and commonly include Raynaud's phenomenon ($>$85\%), puffy hands or sclerodactyly, inflammatory arthritis (non-erosive, resembling SLE), myositis (proximal muscle weakness, elevated creatine kinase), esophageal dysmotility, interstitial lung disease (up to 50\%), pulmonary hypertension (the leading cause of death), serositis, and a characteristic malar rash or discoid lupus-like lesions, with renal involvement being less common and less severe than in SLE alone. Diagnosis is based on the presence of anti-U1-RNP antibodies (high-titer) plus clinical features from at least two of the three component diseases (SLE-like, SSc-like, PM/DM-like), with several classification criteria existing (Alarcón-Segovia, Kahn, Sharp). Management is determined by the dominant clinical manifestation: corticosteroids and disease-modifying antirheumatic drugs (DMARDs) for arthritis and serositis; immunosuppressants (mycophenolate, cyclophosphamide) or rituximab for interstitial lung disease; calcium channel blockers and endothelin receptor antagonists for Raynaud's phenomenon and pulmonary hypertension. Prognosis is variable, with pulmonary hypertension being the leading cause of death.

\subsection{Polymyalgia Rheumatica}
\label{sec:Polymyalgia Rheumatica}
Polymyalgia rheumatica is an inflammatory disorder of the elderly, characterized by pain and stiffness of the shoulder and pelvic girdles, affecting individuals over 50 years of age (peak incidence 70--80 years), with a female predominance of 2--3:1. The etiology is unknown, but genetic associations with HLA-DR4 and environmental triggers have been implicated, with pathophysiology involving local inflammation of proximal joints and extra-articular synovial structures, infiltration of macrophages and CD4+ T cells, elevated IL-6 levels, and increased acute-phase reactants. Clinical features include acute or subacute onset of severe morning stiffness lasting over 30 minutes in the shoulders, neck, and hips, with associated myalgia, constitutional symptoms (low-grade fever, fatigue, weight loss), and difficulty raising arms or rising from a chair, with peripheral arthritis (knees, wrists) occurring in 25--50\%. PMR overlaps significantly with giant cell arteritis (GCA): 15--30\% of PMR patients have concurrent GCA, and 40--60\% of GCA patients have PMR symptoms. Diagnosis is clinical, supported by elevated ESR (often $>$50--100 mm/hr) and CRP, with imaging (ultrasound, MRI) showing synovitis, bursitis, and tenosynovitis. Management is with low-to-moderate dose corticosteroids (prednisone 15--25 mg/day), typically with rapid clinical response within 24--48 hours and slow tapering over 12--18 months. Prognosis is good, with most patients achieving remission within 1--3 years.

\subsection{Relapsing Polychondritis}
\label{sec:Relapsing Polychondritis}
Relapsing polychondritis is a rare, immune-mediated disease characterized by recurrent, progressive inflammation and destruction of cartilaginous structures, particularly in the ears, nose, larynx, tracheobronchial tree, and joints, with an incidence of approximately 0.5--3.5 per million, equal sex distribution, and typically presenting in the fifth to sixth decade. The condition results from an autoimmune attack on type II, IX, and XI collagen and other cartilage matrix components, mediated by autoantibodies, T-cell infiltration, and complement activation, with proteoglycan degradation by matrix metalloproteinases contributing to cartilage destruction. Clinical features include auricular chondritis (the most common presenting feature: painful, swollen, erythematous ears sparing the lobule), nasal chondritis (saddle nose deformity), laryngotracheal involvement (subglottic stenosis, hoarseness, stridor, airway collapse---the most life-threatening manifestation), costochondritis, and non-erosive seronegative arthritis, with ocular involvement (episcleritis, scleritis, uveitis) occurring in up to 60\%, audiovestibular involvement (conductive and sensorineural hearing loss, vertigo) being common, and association with other autoimmune diseases (SLE, vasculitis, myelodysplastic syndrome) in 30\%. Diagnosis is clinical, based on McAdam's criteria or Damiani's criteria, with laboratory findings being non-specific. Management includes corticosteroids (rapid symptomatic relief), dapsone (mild cases), and immunosuppressants (methotrexate, mycophenolate, cyclophosphamide) for severe disease. Prognosis is variable, with respiratory involvement being the leading cause of death.

\subsection{Systemic Sclerosis (Scleroderma)}
\label{sec:Systemic Sclerosis (Scleroderma)}
Systemic sclerosis is a chronic autoimmune connective tissue disease characterized by widespread vasculopathy, immune dysregulation, and progressive fibrosis of the skin and internal organs, with an incidence of approximately 1--2 per 100,000, a female-to-male ratio of 4--6:1, and peaking in the fifth decade. The condition results from excessive collagen deposition by activated fibroblasts, driven by endothelial injury, autoimmune inflammation, and transforming growth factor-beta (TGF-$\beta$) signaling. Clinical features include Raynaud's phenomenon (the most frequent initial manifestation, present in over 95\%), and three clinical variants exist: limited cutaneous SSc (lcSSc, skin thickening distal to elbows/knees, formerly CREST syndrome: Calcinosis cutis, Raynaud's phenomenon, Esophageal dysmotility, Sclerodactyly, Telangiectasias), diffuse cutaneous SSc (dcSSc, skin thickening proximal to elbows/knees and trunk, more rapid and severe internal organ involvement), and systemic sclerosis sine scleroderma (Raynaud's phenomenon and internal organ involvement without clinically apparent skin fibrosis). Pulmonary involvement is the leading cause of death: pulmonary fibrosis (more common in dcSSc) presents with dyspnea, cough, and restrictive lung disease; pulmonary hypertension (more common in lcSSc) presents with exertional dyspnea, right heart failure, and carries a poor prognosis. Scleroderma renal crisis (SRC) is a life-threatening complication of dcSSc (incidence 5--10\%), characterized by malignant hypertension, rapidly progressive renal failure, microangiopathic hemolytic anemia, and thrombocytopenia. Other features include gastroesophageal reflux, gastric antral vascular ectasia (watermelon stomach), telangiectasias, and myositis. Diagnosis is clinical, supported by autoantibodies (anti-centromere in lcSSc, anti-Scl-70/anti-topoisomerase I in dcSSc, anti-RNA polymerase III associated with SRC risk), nailfold capillaroscopy, and organ-specific testing (PFTs, echocardiogram, high-resolution CT chest). Management is organ-specific and multidisciplinary: calcium channel blockers, prostacyclin analogs, and endothelin receptor antagonists for Raynaud's phenomenon and pulmonary hypertension; proton pump inhibitors and prokinetics for GI involvement; ACE inhibitors for scleroderma renal crisis (first-line, which has transformed outcomes---SRC mortality has decreased from 80\% to 15--20\%); mycophenolate, cyclophosphamide, and nintedanib for interstitial lung disease; and methotrexate for skin and joint involvement. Prognosis is variable, with the diffuse cutaneous variant having a worse prognosis, and pulmonary involvement being the leading cause of death.

\subsection{Undifferentiated Connective Tissue Disease}
\label{sec:Undifferentiated Connective Tissue Disease}
Undifferentiated connective tissue disease is a diagnosis applied to patients who have clinical and laboratory features suggestive of a systemic autoimmune disease but do not meet classification criteria for a defined connective tissue disease, estimated to account for 10--50\% of patients presenting to rheumatology clinics with early connective tissue disease symptoms, with a female-to-male ratio of approximately 5:1 and onset typically in the third to fifth decade. Clinical features are mild and often limited to a few domains: the most common are arthralgias or non-erosive arthritis (50--80\%), Raynaud's phenomenon (40--60\%), rash (malar-like or photosensitive), oral ulcers, sicca symptoms, serositis, and constitutional symptoms. Laboratory findings typically include a positive ANA (low-to-moderate titer, often speckled or homogeneous pattern). The hallmark of UCTD is that symptoms remain stable and mild over long-term follow-up, with only 10--30\% of patients evolving into a defined connective tissue disease (most commonly SLE or RA) after 3--5 years. Diagnosis is one of exclusion: the patient must have symptoms compatible with a connective tissue disease for at least 1 year without fulfilling any classification criteria. Management is symptomatic: NSAIDs for arthralgias, hydroxychloroquine for constitutional and cutaneous symptoms, calcium channel blockers for Raynaud's phenomenon, and low-dose corticosteroids for severe symptoms. Prognosis is excellent, with most patients maintaining a stable, mild disease course over decades.

%-------------------------------------------------------------
\section{Variable Vessel Vasculitis}
%-------------------------------------------------------------

%-------------------------------------------------------------

%-------------------------------------------------------------

%-------------------------------------------------------------

%-------------------------------------------------------------

%-------------------------------------------------------------

%-------------------------------------------------------------

\label{sec:Variable Vessel Vasculitis}
%-------------------------------------------------------------%-------------------------------------------------------------

\subsection{Behçet's Disease}
\label{sec:Behçet's Disease (Connective Tissue)}
Behçet's disease is a chronic, relapsing systemic vasculitis of unknown etiology capable of affecting vessels of all sizes, with a predilection for the oral and genital mucosa, eyes, skin, and central nervous system. It is discussed in detail in Chapter~\ref{ch:lymph} (Section~\ref{sec:Behçet's Disease}). Briefly, it is characterized by recurrent oral aphthous ulcers, genital ulcers, uveitis, skin lesions (erythema nodosum, pseudofolliculitis, pathergy reaction), vascular involvement (arterial aneurysms, venous thrombosis), and neurological involvement. It is most prevalent along the Silk Road. Diagnosis is clinical (ICBD criteria). Management is tailored to organ involvement: colchicine for mucocutaneous disease, azathioprine, cyclosporine, or TNF inhibitors for ocular disease, and cyclophosphamide for severe vascular/neurological disease.

\subsection{Cogan's Syndrome}
\label{sec:Cogan's Syndrome}
Cogan's syndrome is a rare, chronic inflammatory disorder characterized by non-syphilitic interstitial keratitis and vestibuloauditory symptoms reminiscent of Ménière's disease/hearing loss, associated with systemic vasculitis, typically affecting young adults (median age 25--40 years), with equal sex distribution. The pathophysiology is poorly understood but likely involves a T-cell-mediated and humoral autoimmune response directed against a common antigen in the cornea and inner ear, with approximately 50--70\% of patients having systemic vasculitis that can involve large, medium, or small vessels. Clinical features include ocular involvement (interstitial keratitis---the hallmark, presenting with eye pain, photophobia, blurred vision, conjunctival injection), vestibuloauditory involvement (acute-onset Ménière-like attacks of vertigo, nausea, vomiting, tinnitus, and fluctuating sensorineural hearing loss that can progress to bilateral profound deafness), and systemic vasculitis (constitutional symptoms, skin nodules or purpura, neuropathy, arthritis). Diagnosis is clinical, based on ocular and vestibuloauditory manifestations in the absence of syphilis or other infection. Management is with corticosteroids (high-dose systemic and topical ophthalmic), with methotrexate, azathioprine, cyclophosphamide, rituximab, and TNF inhibitors used for refractory or severe disease. Prognosis is variable: up to 50\% of patients develop significant visual or hearing impairment.

%-------------------------------------------------------------
\section{Vasculitis (Systemic Vasculitides)}
%-------------------------------------------------------------

%-------------------------------------------------------------

%-------------------------------------------------------------

%-------------------------------------------------------------

%-------------------------------------------------------------

%-------------------------------------------------------------

%-------------------------------------------------------------

\label{sec:Vasculitis (Systemic Vasculitides)}
%-------------------------------------------------------------%-------------------------------------------------------------

The systemic vasculitides are a heterogeneous group of disorders characterized by inflammation and necrosis of blood vessel walls, leading to vessel occlusion, ischemia, and end-organ damage. They are classified according to the size of the predominantly affected vessels (large, medium, small, or variable) as established by the 2012 Chapel Hill Consensus Conference (CHCC).

\subsection{Giant Cell Arteritis}
\label{sec:Giant Cell Arteritis}
Giant cell arteritis is the most common primary systemic vasculitis, characterized by granulomatous inflammation of large and medium-sized arteries, with a predilection for the cranial branches of the aorta, particularly the temporal artery, exclusively affecting individuals over 50 years of age (peak incidence 70--80 years), with a female predominance of 2--3:1, and being more common in Northern European descendants. The condition results from a Th1- and Th17-driven inflammatory response targeting the vessel wall, with CD4+ T-cell infiltration, multinucleated giant cells, intimal hyperplasia, and luminal occlusion, with IL-6 being a key driver. Clinical features include new-onset severe headache (the most common symptom, often temporal), scalp tenderness, jaw claudication (pathognomonic), visual disturbances (amaurosis fugax, diplopia, anterior ischemic optic neuropathy leading to irreversible blindness in 15--20\% if untreated), polymyalgia rheumatica symptoms in 40--60\%, fever, weight loss, and an elevated ESR, with aortic involvement (aortitis, aortic aneurysm, aortic dissection) occurring in 10--20\%. Diagnosis requires a high index of suspicion: ESR is almost always markedly elevated ($>$50--100 mm/hr), CRP is elevated, and temporal artery biopsy (TAB, at least 3--5 cm in length) shows granulomatous inflammation with multinucleated giant cells. Management is with high-dose corticosteroids (prednisone 40--60 mg/day or IV methylprednisolone if vision threatened), with tocilizumab (anti-IL-6 receptor) approved for GCA and used as a steroid-sparing agent. Prognosis is good with prompt treatment, though delayed treatment risks irreversible blindness.

\subsection{Kawasaki Disease}
\label{sec:Kawasaki Disease}
Kawasaki disease is an acute, self-limited, systemic vasculitis of medium-sized arteries, primarily affecting children under 5 years of age, with a predilection for the coronary arteries, with the highest incidence in Japan (300 per 100,000 children $<$5 years), and approximately 20--25 per 100,000 in North America and Europe. The etiology is unknown but suspected to involve an infectious trigger in genetically susceptible individuals, leading to an intense innate and adaptive immune response, with pathophysiology involving systemic inflammation, endothelial activation, and infiltration of the vessel wall by neutrophils, lymphocytes, and macrophages, with consequent arteritis that can lead to coronary artery aneurysm formation and thrombosis. Clinical diagnosis requires fever for at least 5 days plus four of five principal features: bilateral conjunctival injection (non-exudative), oral mucosal changes (strawberry tongue, erythematous lips, pharyngeal erythema), polymorphous rash (usually truncal), extremity changes (erythema and edema of hands/feet in acute phase, periungual desquamation in subacute phase), and cervical lymphadenopathy ($>$1.5 cm, typically unilateral). Complications of coronary arteritis include coronary artery aneurysms (occur in 25\% of untreated KD, reduced to $<$5\% with treatment), myocardial infarction, arrhythmias, and sudden death. Echocardiography is essential at diagnosis and follow-up. Management is intravenous immunoglobulin (IVIG, 2 g/kg single dose) plus high-dose aspirin in the acute phase, reduced to low-dose antiplatelet dose after defervescence. Prognosis is excellent with treatment, with coronary artery aneurysm rates reduced to $<$5\%.

\subsection{Polyarteritis Nodosa}
\label{sec:Polyarteritis Nodosa}
Polyarteritis nodosa is a systemic necrotizing vasculitis affecting medium-sized arteries, characterized by inflammation and fibrinoid necrosis of the vessel wall, with aneurysm formation and thrombosis, sparing small vessels (capillaries, venules) and not associated with ANCA, with an incidence of 2--3 per million, a slight male predominance, and affecting all ages. The condition involves immune complex deposition in the arterial wall (especially in HBV-associated PAN), complement activation, and neutrophil recruitment leading to vessel wall necrosis, with approximately 10--20\% of cases associated with chronic hepatitis B virus (HBV) infection. Clinical features include constitutional symptoms (fever, weight loss, myalgia, arthralgia), mononeuritis multiplex (the most characteristic neurological manifestation), livedo reticularis, testicular pain, renal artery involvement (renal infarction, hypertension, but not glomerulonephritis---distinguishing PAN from microscopic polyangiitis), mesenteric ischemia, and cutaneous manifestations (nodules, ulcers, gangrene). Diagnosis requires biopsy of affected tissue showing necrotizing arteritis with fibrinoid necrosis, or angiography showing microaneurysms and irregular stenoses. Management for HBV-negative PAN is high-dose corticosteroids plus cyclophosphamide for severe disease; HBV-associated PAN requires antiviral therapy plus a shorter course of corticosteroids and plasmapheresis. Prognosis is good with treatment, but severe cases can be fatal.

\subsection{Takayasu Arteritis}
\label{sec:Takayasu Arteritis}
Takayasu arteritis is a chronic, granulomatous large-vessel vasculitis affecting the aorta and its major branches, predominantly in young women, with an incidence of 1--2 per million, a 9:1 female predominance, typically presenting between 10--40 years of age, and being more common in Asia, India, and Latin America. The condition results from Th1-mediated granulomatous inflammation, intimal thickening, fibrosis, and stenosis of large elastic arteries, with skip lesions being characteristic. Clinical features occur in two phases: an early systemic phase (non-specific constitutional symptoms: fever, malaise, weight loss, night sweats, arthralgias) and a late occlusive phase (claudication of extremities, absence or asymmetry of pulses---pulseless disease, blood pressure differences between arms, vascular bruits, syncope, stroke, visual disturbances, renovascular hypertension, and aortic regurgitation). Diagnosis is based on the 1990 ACR criteria, with imaging (conventional angiography, CT angiography, or MR angiography) showing stenosis, occlusion, or aneurysm of the aorta and its branches, and PET-CT showing vessel wall FDG avidity in active disease. Management includes high-dose corticosteroids as first-line therapy, with methotrexate, azathioprine, mycophenolate, TNF inhibitors, and tocilizumab used as steroid-sparing agents or for refractory disease. Prognosis is variable; disease is chronic with relapses requiring long-term immunosuppression.

